\scanpage{387}
Let us note on this occasion that the systems
\[
(31)\qquad
(\alpha,\beta,\unsure{\mu}=1,\varepsilon=1)
\quad\text{with}\quad
\alpha\in\Gamma\deleted{\unsure{\{\pm1\}}}\inserted{\mathbb G_m},\qquad
\beta\in\Gamma\deleted{\unsure{\{\pm1\}}}\inserted{\mathbb G_m}
\]
are in $\unsure{\Gamma|S}\mathcal G_{\rho,\sigma}^{\unsure{\circ|0}}$.
\deleted{and even in $S\mathcal G_{\rho,\sigma}$ if $\alpha,\beta\in\Gamma\mu_2$ ---}
[which are indeed stable under multiplication, \deleted{\unsure{for}} $(\alpha,\beta,\mu,\varepsilon)$, of $\alpha$ and $\beta$ each by a section of $\deleted{\unsure{\mu_2}}\inserted{\mathbb G_m}$, which does not change $\widetilde\alpha,\widetilde\beta\ \ldots$]
they correspond, by the isomorphism $\mathcal G_{\rho,\sigma}\simeq B'_0\times T_\rho\times T_\sigma$, to the systems $(\underline u,\underline a,\underline b)$ with $\underline u=1$, $\underline a\in\Gamma\mathbb G_m$, $\underline b\in\Gamma\mathbb G_m$) and form a \uline{subgroup} (with ordinary multiplication \deleted{\unsure{on the|on the}} \deleted{in $\alpha,\beta$} \inserted{in $\alpha\ \beta$} as group law) --- if $\alpha,\beta$ are sections of $\mu_2$, we find a \unsure{subgroup} \deleted{trivial} ``\unsure{by nature|in action} trivial'' \inserted{$\mu_2\times\mu_2$} of $S\mathcal G_{\rho,\sigma}$.
The action by translation of this subgroup $\simeq\mathbb G_m\times\mathbb G_m$ on $\mathcal G_{\rho,\sigma}^{*}$ is given by multiplication \deleted{in} \inserted{of the} \deleted{the} components in question. This shows that the fact of working with $\underline\alpha,\underline\beta$ defined only ``up to \uline{signs}'' simply amounts to working in the quotient groups of \uline{$\mathcal G_{\rho,\sigma}^{*}(\widehat{\mathbb Z})$} (or of $S\mathcal G_{\rho,\sigma}^{*}(\widehat{\mathbb Z})$) by the subgroup
\[
\inserted{\boxed{\begin{gathered}\Sigma_0\ .\\ \simeq\mathbb Z/2\mathbb Z\times\mathbb Z/2\mathbb Z\end{gathered}}}
\]
\deleted{generated} of the elements $(\alpha,\beta,1,1)$ with $\alpha,\beta\in\{\pm1\}$. This
\margin{386}

% Printed archive footer (not autograph).
\begin{center}
386\\
\texttt{https://grothendieck.umontpellier.fr}\\
Université de Montpellier, UMR CNRS 5149, Case courrier 051, Place Eugène Bataillon\\
34095 Montpellier cedex 5 - France
\end{center}

\scanpage{388}
\mspage{495}

makes precise at the same time the meaning of the expression:
we have a \unsure{homomorphism} of groups \ldots .

Let us now return to the map (23) with arbitrary $f$, and make precise the
composition law of $S_0\mathcal M_{0,3}^{!\sim}$ in terms of
$(\alpha,\beta,f)$ i.e. in terms of the $u_{\alpha,\beta,f}$.
Let us consider the subgroup formed of the $(1,1,f)$,
\deleted{or again} \deleted{of the} which corresponds to the \inserted{pairs}
$(u,f)$ with $u=\operatorname{int}(f^{-1})$, i.e. to the pairs
$(\operatorname{int}(f^{-1}),f)$
--- they form a \unsure{\textquotedblleft}group isomorphic to $\widehat\Pi_{0,3}$,
and we have a $\frac12$ direct product structure:
\[
(31)\qquad
S_0\mathcal M_{0,3}^{!\sim}
\simeq \mathcal M_{0,3}^{!\sim}[0]\cdot\widehat\Pi_{0,3}
\]
where
\[
(32)\qquad
\left\{
\diagram{\begin{tikzcd}[ampersand replacement=\&,column sep=large,row sep=small]
\widehat\Pi_{0,3} \arrow[r] \&
S_0\mathcal M_{0,3}^{!\sim} \& {}\\
f \arrow[r,mapsto] \&
(\deleted{\unsure{f,}}\operatorname{int}(f),f)
\arrow[r,leftrightarrow] \& (1,1,f^{-1}),
\end{tikzcd}}
\right.
\]
the operation of $\mathcal M_{0,3}^{!\sim}[0]$ on $\widehat\Pi_{0,3}$
\deleted{\unsure{being}} the tautological operation
(deduced from the inclusion
$\mathcal M_{0,3}^{!\sim}[0]\subset\mathcal M_{0,3}^{!\sim}
\subset\underline{\operatorname{Aut}}_{\mathrm{lac}}
(\widehat\Pi_{0,3})$).

The \inserted{invariant} subgroup $\widehat{\mathbb Z}$ of
$S_0\mathcal M_{0,3}^{!\sim}$ is given by the pairs
$\inserted{(u,f)=}(1,\ell_0^n)$ $(n\in\widehat{\mathbb Z})$
which we also put in the form $(\alpha,\beta,f)$
where $\alpha,\beta$ are such that
$u_{\alpha,\beta}=fu=\operatorname{int}(\ell_0^n)$,
hence $\alpha=\beta=\ell_0^n$, i.e.
\[
(33)\qquad
\widehat{\mathbb Z}\longrightarrow S_0\mathcal M_{0,3}^{!\sim}
\quad\text{is identified with}\quad
n\longmapsto
(\deleted{\unsure{\alpha=}}\ell_0^n,\beta=\ell_0^n,f=\ell_0^n)
\]
\margin{387}

% Printed archive footer (not autograph).
\begin{center}
387\\
\texttt{https://grothendieck.umontpellier.fr}\\
Université de Montpellier, UMR CNRS 5149, Case courrier 051, Place Eugène Bataillon\\
34095 Montpellier cedex 5 - France
\end{center}

\scanpage{389}

We \unsure{will have}
\[
\begin{aligned}
(f'^{-1}u_{\alpha',\beta'})(f^{-1}u_{\alpha,\beta})
&=\deleted{\unsure{f'^{-1}f^{-1}(f u_{\alpha',\beta'}f^{-1})}}\\
&\inserted{\underbrace{f'^{-1}\unsure{u_{\alpha',\beta'}}(f^{-1})}
_{\unsure{f''^{-1}}}
(u_{\alpha',\beta'}u_{\alpha,\beta})}
\end{aligned}
\]
\unsure{i.e.} the multiplication in $S_0\mathcal M_{0,3}^{!\sim}$
is given, in terms of the components
$(\alpha',\beta',f')(\alpha,\beta,f)$,
\unsure{for|by} \deleted{\unsure{g''=}}$(\alpha'',\beta'',f'')$, where
\[
(34)\quad
\left\{
\begin{aligned}
u_{\alpha'',\beta''}
&=u_{\alpha',\beta'}u_{\alpha,\beta}
\quad\text{i.e.}\quad
\left\{
\begin{aligned}
\alpha''&=u_{\alpha',\beta'}(\alpha)\alpha'\\
\beta''&=u_{\alpha',\beta'}(\beta)\beta'
\end{aligned}
\right.,
\quad\text{and}\\
f''&=u_{\alpha',\beta'}(f)f'
\end{aligned}
\right.
\]

Returning now to the definition of $S\mathcal M_{\rho\sigma}$
(1), we are going to introduce there the \unsure{corresponding} multiplication,
which amounts to defining $S\mathcal M_{\rho\sigma}$ as
a $\frac12$ direct product
\[
\begin{aligned}
(35)\qquad S\mathcal M_{\rho,\sigma}
&\overset{\mathrm{dfn}}{=}
S\mathcal G_{\rho,\sigma}^{*}\cdot SG\\
&\qquad(\text{product }\tfrac12\text{ direct})
\end{aligned}
\]
\uline{by the operation of} $S\mathcal G_{\rho,\sigma}^{*}$
\uline{on} $G$
\uline{via the homomorphism} $\Theta$ (page~\unsure{489'|488'}),
and \unsure{by} \unsure{defining|will define}
\unsure{the|the} \deleted{map}\inserted{morphism}
\[
(36)\qquad
\left|
\diagram{\begin{tikzcd}[ampersand replacement=\&,column sep=large,row sep=small]
S\mathcal M_{\rho,\sigma}\arrow[r] \&
S\mathcal G_{\rho,\sigma}^{*}\times_S G\\
\deleted{\unsure{(\alpha,\beta,\mu,\varepsilon)|(\alpha,\beta)}}
\inserted{U}\cdot f\arrow[r,mapsto] \& (U,f^{-1})
\end{tikzcd}}
\right.
\]
\margin{Watch out for the passage from $f\underset{\leftarrow}{\ \unsure{\text{to}}\ }f^{-1}$!}
where $U=(\alpha,\beta,\mu,\varepsilon)$ is in
$\Gamma S\mathcal G_{\rho,\sigma}$.
\margin{388}

% Printed archive footer (not autograph).
\begin{center}
388\\
\texttt{https://grothendieck.umontpellier.fr}\\
Université de Montpellier, UMR CNRS 5149, Case courrier 051, Place Eugène Bataillon\\
34095 Montpellier cedex 5 - France
\end{center}

\scanpage{390}
\mspage{496}

We have thus done what was needed, in the case
$S=\operatorname{Spec}(\widehat{\mathbb Z})$ etc, to have a
\uline{homomorphism of groups}
\[
(37)\qquad
\diagram{\begin{tikzcd}[ampersand replacement=\&,column sep=large]
S_0\mathcal M_{0,3}^{!\sim}\arrow[r] \&
S\mathcal M_{\rho,\sigma}(\widehat{\mathbb Z})/
\unsure{\widetilde{\Sigma}}\deleted{\unsure{(\widehat{\mathbb Z})}}
\end{tikzcd}}
\]
where $\Sigma_0\simeq\mathbb Z_{/2}\times\mathbb Z_{/2}$
is the \unsure{subgroup} \deleted{\unsure{\textquotedblleft powerless\textquotedblright}}
of $S\mathcal G_{\rho,\sigma}^{*}$ introduced on page~494',
which operates trivially on $G$ hence is invariant in
$S\mathcal M_{\rho,\sigma}$, hence in
$S\mathcal M_{\rho,\sigma}(\widehat{\mathbb Z})$, and where
\[
(38)\quad
\widetilde{\Sigma}=\Sigma_0\times
\underbrace{\{\pm1\}}_{
\substack{
\text{this factor identified with the center of }SG(\unsure{\widehat{\mathbb Z}})\\
=\operatorname{Sl}(2,\widehat{\mathbb Z}),\
\text{contained in }\operatorname{Gl}(2,\widehat{\mathbb Z})\\
=G(\widehat{\mathbb Z})
}}
\simeq
\mathbb Z/2\mathbb Z\times\mathbb Z/2\mathbb Z\times\mathbb Z/2\mathbb Z
\]
is the corresponding normal group in the $\frac12$ direct product
--- the fact of passing to the quotient by $\Sigma$ expresses
the three \unsure{ambiguities|ambiguities} of \uline{signs},
on $\alpha$, on $\beta$ and on $f$.

Let us now consider the \unsure{homomorphism}
\[
(39)\quad
\left|
\diagram{\begin{tikzcd}[ampersand replacement=\&,column sep=large,row sep=small]
\mathbb G_a\arrow[r] \&
\unsure{S_0|S}\mathcal M_{\rho,\sigma}^{\deleted{\unsure{!\sim}}}\\
\lambda\arrow[r,mapsto] \&
(\alpha\unsure{=}\varepsilon_0(\unsure{\lambda|-\lambda}),
\beta\unsure{=}\varepsilon_0(-\lambda),
\mu=1,\varepsilon=1,f=\varepsilon_0(\lambda))
\end{tikzcd}}
\right.
\]
it is a homomorphism of groups, and even makes
$\mathbb G_a$ a \uline{normal} subgroup.
To \inserted{verify} this the most convenient thing is to identify
$\unsure{S_0|S}\mathcal M_{\rho,\sigma}^{!\sim}$
with the $\frac12$-direct product (35),
\margin{389}

% Printed archive footer (not autograph).
\begin{center}
389\\
\texttt{https://grothendieck.umontpellier.fr}\\
Université de Montpellier, UMR CNRS 5149, Case courrier 051, Place Eugène Bataillon\\
34095 Montpellier cedex 5 - France
\end{center}

\scanpage{391}

where we identify $S\mathcal G_{\rho,\sigma}^{*}$ with a subgroup
of $B'_0\times T_\rho\times N'_\sigma$, operating on $G$
by inner automorphisms \uline{via} $B'_0$.

The preceding morphism thus corresponds to
\[
(40)\qquad
\diagram{\begin{tikzcd}[ampersand replacement=\&,column sep=small,row sep=small]
\lambda\arrow[r,mapsto,"i"]\arrow[d,phantom,"\in"{rotate=90}] \&
{(\lambda,1,1,} \&
\lambda^{-1}\arrow[d,phantom,"\in"{rotate=90}] \&
{)=\deleted{\lambda\cdot\unsure{j}(\lambda)}}\\
\Gamma B'_0 \& {} \& \Gamma G \& {}
\end{tikzcd}}
\]
\margin{\deleted{\diagram{\begin{tikzcd}[ampersand replacement=\&,column sep=small]
\unsure{j}:SG\arrow[r] \& S\mathcal M_{\rho,\sigma}
\end{tikzcd}} is} the canonical homomorphism of groups}

\margin{\deleted{\unsure{where $U_0$}}}
\deleted{Let on the other hand
$U=(\underbrace{\underline u,\underline{\unsure{a|\rho}},
\underline{\unsure{b|\beta}}}_{U_0},\unsure{\varepsilon}f)$
\inserted{$=U_0\cdot f$},
be an arbitrary element of the $\frac12$ direct product,
and let us compute $U i(\lambda)U^{-1}$ ---
\[
U i(\lambda)U^{-1}=U_0 f i(\lambda)f^{-1}U_0^{-1}=
\]
}

We identify $S\mathcal G_{\rho,\sigma}^{*}$ with a subgroup of
$\mathcal G_{\rho,\sigma}^{*}$, and likewise
the $\frac12$ direct product $S\mathcal M_{\rho,\sigma}^{\unsure{*}}$
with a subgroup of the $\frac12$ direct product
\[
(41)\qquad
\deleted{\unsure{S\mathcal M_{\rho,\sigma}^{*}}}
\inserted{\widetilde S\mathcal M_{\rho,\sigma}}
=\mathcal G_{\rho,\sigma}^{*}\cdot G
\simeq(B'_0\times T_\rho\times N'_\sigma)\cdot G,
\]
and we are going to see that the image of the \unsure{monomorphism} $i$ (40) in
$\mathcal G_{\rho,\sigma}^{*}\cdot G$ is even
normal in this larger group.
So as not to get lost in the notation,
let us denote by
\[
(42)\qquad i_B,\ i_\rho,\ i_\sigma,\ i_G
\]
the morphisms \unsure{of inclusions|of inclusion}
\[
\left|
\diagram{\begin{tikzcd}[ampersand replacement=\&,column sep=large,row sep=small]
B'_0\arrow[r] \& \deleted{\unsure{\mathcal G_{\rho,\sigma}^{*}}}\\
T_\rho\arrow[r] \& {}\\
\unsure{N'_\sigma|N''_\sigma}\arrow[r] \& {}\\
G\arrow[r] \& {}
\end{tikzcd}}
\right\}\ 
\unsure{i'}
\begin{gathered}
\inserted{\widetilde S\mathcal M_{\rho,\sigma}}\\
\deleted{\unsure{\mathcal G_{\rho,\sigma}^{*}|\mathcal M_{\rho,\sigma}}}\cdot G
\end{gathered}
\]
\margin{390}

% Printed archive footer (not autograph).
\begin{center}
390\\
\texttt{https://grothendieck.umontpellier.fr}\\
Université de Montpellier, UMR CNRS 5149, Case courrier 051, Place Eugène Bataillon\\
34095 Montpellier cedex 5 - France
\end{center}